\documentclass[12pt]{article}

%%%%%%%%%%%%%%%%%%%%%%%%%%%%%%%%%%%%%%%%%%%%%%%%%%%%%%%%%%%%%%%%%%
%                     PREAMBLE / PACKAGES
%%%%%%%%%%%%%%%%%%%%%%%%%%%%%%%%%%%%%%%%%%%%%%%%%%%%%%%%%%%%%%%%%%

% ----------------------------------------------------------------
% Font encoding and page geometry
% ----------------------------------------------------------------
\usepackage[T1]{fontenc}
\usepackage[margin=1in]{geometry}

% ----------------------------------------------------------------
% Line and paragraph spacing
% ----------------------------------------------------------------
\usepackage{setspace}
\usepackage{parskip}

% ----------------------------------------------------------------
% Section and table-of-contents formatting
% ----------------------------------------------------------------
\usepackage{titlesec}
\usepackage{tocloft}

% ----------------------------------------------------------------
% Graphics, figures, and page layouts
% ----------------------------------------------------------------
\usepackage{graphicx}
\usepackage{animate} % frame-by-frame animations in supported PDF viewers
\usepackage[font=small]{caption}
\usepackage{subcaption}
\usepackage{float}
\usepackage{placeins}
\usepackage{wrapfig}
\usepackage{multicol}

% ----------------------------------------------------------------
% Mathematics
% ----------------------------------------------------------------
\usepackage{amsmath}
\usepackage{amssymb}
\usepackage{amsthm}
\usepackage{mathrsfs}
\usepackage{cancel}
\usepackage{gensymb}

% ----------------------------------------------------------------
% Tables and lists
% ----------------------------------------------------------------
\usepackage{booktabs}
\usepackage{tabularx}
\usepackage{multirow}
\usepackage{enumitem}

% ----------------------------------------------------------------
% Colors and highlighting
% ----------------------------------------------------------------
\usepackage{xcolor}
\usepackage{soul}
\sethlcolor{yellow}

% ----------------------------------------------------------------
% Headers and footers
% ----------------------------------------------------------------
\usepackage{fancyhdr}

% ----------------------------------------------------------------
% Source-code examples
% ----------------------------------------------------------------
\usepackage{listings}

% ----------------------------------------------------------------
% Citations and hyperlinks
% ----------------------------------------------------------------
\usepackage[numbers]{natbib}
\usepackage{hyperref}

% ----------------------------------------------------------------
% Default line spacing
% ----------------------------------------------------------------
\onehalfspacing
% ------------------------------------------------------------
% Document Information
% ------------------------------------------------------------
\title{Lab 5}
\author{Chuan Wei Heng}
\date{\today}

% ------------------------------------------------------------
% Document
% ------------------------------------------------------------
\begin{document}

\maketitle

For today’s lab, we explore the effectiveness of the 
Bisection and Newton’s methods and use a 
guaranteed-convergence criterion for Newton’s method to 
construct a hybrid root-finding method.

Let us first discuss this criterion derived for 
guaranteed-convergence. We first recall that the Newton’s
method algorithm is as follows:

\[
x_{n+1}=x_{n}-\frac{f(x_n)}{f'(x_n)}.
\]

We also recall the fixed point convergence theorem discussed
previously which states that $\forall x_0 \in (x^*-\varepsilon,
x^*+\varepsilon)$, the iteration will converge to the root 
$x^*$ if $|f'(x)|<1.$ In other words, if the derivative is
bounded by 1 in some $\varepsilon$ neighborhood centered 
about $x^*.$

With that in mind, we treat the Newton’s method iteration as
a fixed point iteration, by treating $x_{n+1}$ as $g(x)$, we 
can obtain the following:

\begin{align}
g'(x) 
& = 1-\frac{f'(x)^2-f(x)f''(x)}{f'(x)^2},
\\
& = \frac{f(x)f''(x)}{f'(x)}.
\end{align}

With this, we obtain a sufficient criterion for the 
guaranteed convergence of Newton’s method. We first introduce
the function of interest. Then, we discuss the convergence 
behavior of the Bisection and Newton’s methods before 
proposing a hybrid approach.

\section*{Function of Interest}

The function we study is $f(x)=e^{x^2+7x-30}-1$. We provide a plot below for visualization.

\begin{figure}[H]

    \centering

    \includegraphics[width=0.75\linewidth]{images/plotting_the_function.png}

    \caption{Plot of the function $f(x)=e^{x^2+7x-30}-1$.}

    \label{fig:single_example}

\end{figure}

We now present the various numerical root-finding methods 
used to determine the roots of this function.

\section*{Bisection Method}

\begin{figure}[H]

    \centering

    \includegraphics[width=0.75\linewidth]{images/bisection.png}


    \label{fig:single_example}

\end{figure}

\section*{Newton's Method}

\begin{figure}[H]

    \centering

    \includegraphics[width=0.75\linewidth]{images/Newton.png}


    \label{fig:single_example}

\end{figure}

\section*{Proposed Hybrid Method}
We now propose a hybrid method that should perform more 
efficiently than the two previously presented methods. This 
is primarily motivated by the fact that we cannot expect the 
user to always have a good initial guess for the location of 
a root, particularly since not every user of a root-finding 
method will be a numerical analyst.

We therefore consider the main strengths and weaknesses of 
the Bisection and Newton's methods. The Bisection method has 
a wide range of convergence but converges relatively slowly, 
while Newton's method converges rapidly when supplied with a 
suitable initial guess but can be highly sensitive to that 
choice.

Thus, we seek to combine the strengths of the two methods 
while mitigating their respective weaknesses. The proposed 
schematic is as follows: utilize the Bisection method to 
find the basin of convergence for the Newton iteration, and 
switch to Newton's method from there. We present the 
resulting convergence profile bellow.

\begin{figure}[H]

    \centering

    \includegraphics[width=0.75\linewidth]{images/hybrid.png}

    \label{fig:single_example}

\end{figure}

As we can see, the number of iterations is $3-4$ times faster 
than previous methods. Very cool!

For details of the code, we turn the reader towards the git 
repository by \textit{chuanweiheng-cell} entitled APPM4600.
Thank you!

\end{document}